\section{Experimental Study}
\label{sec:experiments}

\subsection{Scenario tree}
\label{sec:scentree}
We use a two-stage design: one-factor-at-a-time (OFAT) screening to rank the
levers, followed by a $2^k$ factorial on the highest-impact levers, then a
demand sweep for 2030 and the Industry~4.0 scrap scenario. The OFAT branch is
ordered to test the analytically predicted constraints
(Sec.~\ref{sec:capacity}) first.

\begin{figure}[H]
\centering
\begin{forest}
  for tree={
    grow'=0, draw, rounded corners, font=\footnotesize,
    edge={-Stealth}, l sep=10mm, s sep=2mm, anchor=west,
    child anchor=west, parent anchor=east, align=left}
  [Baseline\\(bottleneck ID)
    [OFAT screening
      [+1 Drying / faster clean]
      [+1 Cleaning team]
      [Resequence (group by product)]
      [+1 Packaging line]
    ]
    [$2^k$ factorial\\(top levers)
      [Best combination]
    ]
    [2030 demand sweep\\($\times$1.25 / $\times$1.5 / $\times$2)
      [Max sustainable throughput]
    ]
    [Industry 4.0\\(scrap 10\%$\rightarrow$4\%)]
  ]
\end{forest}
\caption{Scenario tree for the experimental study.}
\label{fig:scentree}
\end{figure}

\subsection{Factorial design}
\label{sec:factorial}
OFAT screening ranks the three capacity stations above resequencing, so the
$2^3$ factorial is run on drying, cleaning-team and packaging capacity. A full
factorial (all eight corners, 20 replications each) is used rather than OFAT so
that interaction effects are estimable: OFAT would miss them and can misread
random variation between two runs as an effect.

\begin{table}[H]
\centering
\caption{$2^3$ factor settings (low $=-$, high $=+$).}
\label{tab:factorial}
\begin{tabularx}{\textwidth}{l X X}
\toprule
Factor & Low ($-$) & High ($+$) \\
\midrule
A: Drying capacity        & 1 dryer  & 2 dryers \\
B: Cleaning-team capacity & 1        & 2 \\
C: Packaging capacity     & 3 lines  & 4 lines \\
\bottomrule
\end{tabularx}
\end{table}

Each effect is the difference between the four high-corner means and the four
low-corner means; its standard error is pooled from the within-corner variance
across all eight corners ($\mathrm{SE}=346$ bottles, $\mathrm{df}=152$), so each
effect carries a significance test.

\begin{table}[H]
\centering
\caption{$2^3$ effect estimates on weekly throughput (bottles). Sorted by
magnitude. $|t|>1.99$ is significant at \SI{95}{\percent}.}
\label{tab:doe}
\begin{tabularx}{\textwidth}{X r r c}
\toprule
Effect & Estimate (bottles) & $t$ & Significant \\
\midrule
\textbf{BC: clean $\times$ packaging} & \textbf{$+700$} & \textbf{2.02} & \textbf{yes} \\
AB                         & $-478$ & $-1.38$ & no \\
C: packaging line          & $-473$ & $-1.37$ & no \\
AC                         & $+467$ & 1.35  & no \\
A: dryer                   & $+425$ & 1.23  & no \\
ABC                        & $-308$ & $-0.89$ & no \\
B: cleaning team           & $+80$  & 0.23  & no \\
\bottomrule
\end{tabularx}
\end{table}

\begin{figure}[H]
\centering
\includegraphics[width=0.55\textwidth]{fig_doe.pdf}
\caption{Pareto of $2^3$ effects on throughput. No main effect is significant;
only the cleaning-team$\times$packaging interaction (BC) marginally crosses the
\SI{95}{\percent} threshold.}
\label{fig:doe}
\end{figure}

Table~\ref{tab:doe} and the Pareto in Figure~\ref{fig:doe} show that \emph{no
main effect} is significant: adding a dryer, a cleaning-team member or a
packaging line each sits inside the noise band and cannot be distinguished from
zero. The only effect to cross the threshold is the cleaning-team$\times$packaging
interaction (BC, $+700$ bottles, $t=2.02$), and only marginally. With seven
effects tested at the \SI{5}{\percent} level, one borderline crossing is the
expected rate of a false positive, so we read it as a multiple-comparison
artefact rather than a real synergy. The central result, which the OFAT design
would have hidden, is that no single capacity lever and no robust combination
changes \emph{shipped output}, because output is fixed by the 50-batch plan and
per-batch yield (Sec.~\ref{sec:demand40}), not by station capacity. Adding
capacity moves the bottleneck (Sec.~\ref{sec:btshift}) but not the throughput.

\subsection{Results: throughput across scenarios}
\label{sec:expresults}
Figure~\ref{fig:expresults} compares weekly throughput across the scenario set
with \SI{95}{\percent} confidence intervals. Almost every bar overlaps the
baseline: neither the capacity levers nor the demand sweep moves shipped output.
The one clear exception is Industry~4.0 (press scrap $10\rightarrow\SI{4}{\percent}$),
which lifts throughput by 3{,}551 bottles ($+\SI{5.8}{\percent}$,
95\% CI $[2{,}206,\,4{,}896]$), the only intervention whose interval
clears the baseline by a wide margin.

\begin{figure}[H]
\centering
\includegraphics[width=0.72\textwidth]{fig_throughput.pdf}
\caption{Weekly throughput by scenario with \SI{95}{\percent} CI error bars.
Industry~4.0 (green) is the only scenario that separates from the baseline.}
\label{fig:expresults}
\end{figure}

\paragraph{Queue discipline.} Switching the dryer queue from FIFO to EDD or SPT
(\texttt{queueStrategy}~$=1,2$) leaves throughput at or below the baseline
(EDD $-\SI{1.2}{\percent}$, SPT $-\SI{2.4}{\percent}$); neither improves output,
consistent with throughput being plan-capped. Drying times are near-homogeneous
($\sigma\approx\SI{0.39}{\hour}$, Sec.~\ref{sec:input}), so reordering a queue of
similar jobs cannot meaningfully change flow, so sequencing is not a lever here.

\subsection{Bottleneck shift across scenarios}
\label{sec:btshift}
We report per-station utilisation, not only throughput, to show \emph{which}
resource binds as capacity is added. Table~\ref{tab:btshift} and
Figure~\ref{fig:util} trace the bottleneck as it moves.

\begin{table}[H]
\centering
\caption{Mean utilisation (\%) by station. Bold $=$ binding constraint. The
bottleneck moves from drying to the cleaning team, while throughput stays flat.}
\label{tab:btshift}
\begin{tabularx}{\textwidth}{X c c c c}
\toprule
Scenario & Drying & Clean team & Packaging & Throughput \\
\midrule
Baseline              & \textbf{94.9} & 81.0 & 53.0 & 61{,}761 \\
$+1$ dryer            & 77.7 & \textbf{95.0} & 73.0 & 61{,}890 \\
$+1$ dryer, $+$clean team & 60.9 & \textbf{93.1} & 92.8 & 61{,}099 \\
Best combination ($+$dryer, $+$clean, $+$pack) & 76.3 & \textbf{93.0} & 77.1 & 61{,}485 \\
\bottomrule
\end{tabularx}
\end{table}

\begin{figure}[H]
\centering
\includegraphics[width=0.72\textwidth]{fig_utilisation.pdf}
\caption{Resource utilisation by scenario. Adding a dryer pushes drying off the
\SI{90}{\percent} line but drives the cleaning team onto it, so the constraint
shifts rather than clears.}
\label{fig:util}
\end{figure}

Drying binds at baseline (\SI{94.9}{\percent}). A second dryer relieves it
(\SI{77.7}{\percent}) but immediately loads the cleaning team to
\SI{95.0}{\percent}, because each dryer carries a 2-hour clean that draws on the
same shared team. Adding a cleaning-team member then splits the load between the
team and packaging, both near \SI{93}{\percent}. Throughput never follows: it
holds at 61{,}000--62{,}000 bottles across all four configurations. This is
the bottleneck-shifting pattern: capacity relocates the
constraint but cannot raise output, because output is plan- and yield-limited,
not capacity-limited.

\subsection{2030 demand and Industry 4.0}
\label{sec:demand40}
The \texttt{demandMultiplier} parameter scales the \emph{interarrival times} of
the 50-batch plan, so a higher multiplier releases the same batches faster rather
than adding batches (the run still terminates at the 50\textsuperscript{th}
shipment). The sweep therefore tests whether the line can absorb a compressed
2030 schedule while still honouring the one-day release rule, not whether it can
ship more units. Table~\ref{tab:demand} reports the baseline under
$\times1.25$, $\times1.5$ and $\times2.0$.

\begin{table}[H]
\centering
\caption{Baseline under a compressed 2030 schedule. Throughput is plan-fixed;
the load is absorbed as congestion (flow time, WIP) with zero release
violations.}
\label{tab:demand}
\begin{tabularx}{\textwidth}{X c c c c}
\toprule
Demand & Throughput & Mean flow time (h) & Avg WIP & Release viol. \\
\midrule
$\times$1.0 (baseline) & 61{,}761 & 182.3 & 19.8 & 0 \\
$\times$1.25           & 60{,}542 & 189.0 & 20.6 & 0 \\
$\times$1.5            & 59{,}970 & 193.9 & 21.1 & 0 \\
$\times$2.0            & 60{,}682 & 203.1 & 22.2 & 0 \\
\bottomrule
\end{tabularx}
\end{table}

Throughput stays flat at $\approx61{,}000$ bottles across the whole sweep,
confirming that shipped volume is set by the plan and yield, not by demand. What
does change is congestion: at $\times2.0$ mean flow time rises \SI{11}{\percent}
(\SI{182}{\hour}$\rightarrow$\SI{203}{\hour}) and average WIP rises
\SI{12}{\percent} (19.8$\rightarrow$22.2). Release violations stay at zero
throughout, so the line absorbs a doubled-pace schedule without breaching the
one-day release constraint. Drying remains the binding
resource (\SI{95}{\percent}); it is the station that would constrain any genuine
\emph{volume} growth, since lifting the batch count would push it past
saturation first.

\paragraph{Industry 4.0: the only lever that raises output.}
Reducing press scrap from \SI{10}{\percent} to \SI{4}{\percent} raises good
tablets per batch, and because output is bounded by batches $\times$ yield, this
is the one intervention that lifts \emph{shipped} throughput: 65{,}313 vs.\
61{,}761 bottles, a gain of \textbf{3{,}551 bottles per run}
($+\SI{5.8}{\percent}$, 95\% CI $[2{,}206,\,4{,}896]$). It is the only lever
that significantly raises shipped output, where no capacity lever does, and it
requires no added line capacity.

\subsection{Managerial value: benefit vs.\ effort}
\label{sec:value}
\begin{table}[H]
\centering
\caption{Each lever's throughput gain against implementation effort. ``n.s.''
$=$ not statistically significant at \SI{95}{\percent} (within the noise band).}
\label{tab:value}
\begin{tabularx}{\textwidth}{X c X}
\toprule
Lever & $\Delta$ throughput & Implementation effort / cost \\
\midrule
Resequence (group by product) & $-0.7\%$ (n.s.) & Free: scheduling change only \\
$+1$ cleaning team            & $+0.1\%$ (n.s.) & Low: additional shift / hire \\
$+1$ dryer                    & $+0.7\%$ (n.s.) & Medium: equipment, moderate lead time \\
$+1$ packaging line           & $-0.8\%$ (n.s.) & High: capex, long lead time (the 2030 plan defers this) \\
\textbf{Industry 4.0 (scrap $10\rightarrow\SI{4}{\percent}$)} & \textbf{$+5.8\%$} & Low: one-off press-sensor retrofit, no recurring cost \\
\bottomrule
\end{tabularx}
\end{table}

Only one lever moves throughput at all: the press-sensor retrofit ($+\SI{5.8}{\percent}$),
at a fraction of the capital cost of any line investment. The capacity levers that
dominate the 2030 plan's intuition (more dryers, more cleaning staff, more
packaging, resequencing) return nothing measurable on output; they buy flow-time
and bottleneck relief, not volume.